% !PRINTED 11
\weberpage{11}
In die Mannigfaltigkeit der rationalen Zahlen ordnen sich die
natürlichen Zahlen selbst mit ein, wenn man unter einer natürlichen
Zahl $m$ das Verhältniss $m:1$ versteht.

Wir kehren jetzt zu irgend einer messbaren Mannigfaltigkeit
$\mathfrak M$ zurück und nehmen aus ihr irgend zwei Elemente $a$ und
$b$. Wählt man, was immer möglich ist, zwei natürliche Zahlen $m,n$,
so dass $na>mb$, so heisst das Verhältniss $a:b$ grösser als das
rationale Verhältniss $m:n$ oder
\[
 {a\over b}>{m\over n},
\]
und wenn $m:n>p:q$, so ist auch $a:b>p:q$. Ebenso folgt, wenn
$n'a<m'b$ ist,
\[
 {a\over b}<{m'\over n'}.
\]

Ist $a:b>m:n$, so kann man eine und folglich auch beliebig viele
rationale Zahlen $m_1:n_1$ finden, so dass
\[
 {a\over b}>{m_1\over n_1}>{m\over n}
\]
ist, d.~h. man kann \gesperrt{zwischen} $a:b$ und $m:n$ beliebig viele
rationale Verhältnisse einschalten. Um dies zu zeigen, wähle man eine
beliebige ganze Zahl $k$ und bestimme die ganze Zahl $h$ so, dass
$h(na-mb)>kb$ wird, was immer möglich ist; dann ist
\[
 {a\over b}>{hm+k\over hn}>{m\over n}.
\]
Und ebenso folgt, wenn $n'a<m'b$, $h'(m'b-n'a)>k'a$ ist,
\[
 {a\over b}<{h'm'\over h'n'+k'}<{m'\over n'}.
\]

Wenn nun $a:b$ und $\alpha:\beta$ irgend zwei Verhältnisse sind, die
wir der Kürze wegen auch mit $e,\epsilon$ bezeichnen wollen, deren
Elemente derselben oder auch verschiedenen Mannigfaltigkeiten
angehören, so sind zwei Fälle möglich: 1) Es liegt kein rationales
Verhältniss $\mu$ zwischen $e$ und $\epsilon$ oder 2) es liegt ein
rationales Verhältniss zwischen $e$ und $\epsilon$.

Im Falle 1) heissen die beiden Verhältnisse $e$ und $\epsilon$
\gesperrt{einander gleich}, und man sieht, dass zwei Verhältnisse, die
einem dritten gleich sind, auch unter einander gleich sind; denn ist
$e<\mu<\epsilon$, so ist jedes andere Verhältniss $e'$ entweder kleiner
%CONT>>
